\subsection{Sequence diagram cho các luồng chính}
\label{sec:sequence-flows}

Ba luồng dưới đây được chọn theo tiêu chí: rủi ro kiến trúc cao nhất, có quyết định thiết kế cần bảo vệ, và trace được về một use case trọng tâm (Bảng~\ref{tab:uc-relations}). Mỗi sơ đồ chỉ mô tả một kịch bản cụ thể; các nhánh điều kiện của cùng use case được liệt kê ở luồng thay thế của đặc tả tương ứng thay vì vẽ nhiều khung \texttt{alt} trong cùng hình.

\subsubsection{Xung đột âm thanh khi mở luồng tiếng thứ hai (X0)}
\label{sec:seq-conflict}

Hình~\ref{fig:seq-conflict} mô tả kịch bản người dùng mở một bài nhạc trong khi video đang phát. Session Coordinator kiểm tra quy tắc đã lưu; khi chưa có quy tắc, người dùng nhận popup hai lựa chọn kèm ô ``Không hỏi lại''. Lựa chọn ``Tắt tiếng video'' giữ video tiếp tục hiển thị nhưng đặt âm lượng bằng không, đồng thời nhạc trở thành nguồn tiếng chính; lựa chọn ``Dừng video'' dừng video và lưu vị trí. Nếu ô nhớ được tích, quy tắc được ghi lại và các lần sau áp dụng im lặng, đặt lại bằng mục ``Luôn hỏi'' trong Cài đặt.

\begin{figure}[H]
\centering
\begin{tikzpicture}[
  font=\tiny,
  ll/.style={draw=black!70, thick, rounded corners=2pt, fill=black!5, minimum width=2.15cm, minimum height=0.5cm, align=center},
  msg/.style={-{Latex[width=2mm, length=2.4mm]}, thick},
  ret/.style={-{Latex[width=2mm, length=2.4mm]}, thick, dashed},
  note/.style={draw=black!55, fill=black!3, rounded corners=2pt, inner sep=2.5pt, align=center}]
  \node[ll] (u)  at (1.1,6.45) {\bfseries Người dùng};
  \node[ll] (ui) at (3.9,6.45) {Màn Nhạc (UI)};
  \node[ll] (sp) at (6.8,6.45) {Session\\Coordinator};
  \node[ll] (mc) at (9.7,6.45) {Media\\Coordinator};
  \node[ll] (pl) at (12.55,6.45) {2 Player\\(nhạc $\cdot$ video)};
  \foreach \n in {u,ui,sp,mc,pl} \draw[dashed, black!50] (\n.south) -- ++(0,-4.6);
  % Thanh kích hoạt
  \foreach \bx/\ya/\yb in {3.9/5.4/5.9, 6.8/1.75/5.4, 9.7/2.45/3.5, 12.55/2.45/3.0}{
    \draw[fill=white, draw=black!60, thick] (\bx-0.07,\ya) rectangle (\bx+0.07,\yb);}
  \draw[msg] (1.1,5.9) -- (3.9,5.9) node[midway, above]{1. Mở bài nhạc (video đang phát)};
  \draw[msg] (3.9,5.4) -- (6.8,5.4) node[midway, above]{2. Xin quyết định luồng tiếng};
  \node[note] at (8.6,5.0) {chưa có quy tắc đã lưu};
  \draw[dashed, black!55] (6.8,5.28) -- (7.6,5.03);
  \draw[ret] (6.8,4.5) -- (1.1,4.5) node[midway, above]{3. Popup 2 lựa chọn + ô tick (X0)};
  \draw[msg] (1.1,4.0) -- (6.8,4.0) node[midway, above]{4. Chọn ``Tắt tiếng video''};
  \draw[msg] (6.8,3.5) -- (9.7,3.5) node[midway, above]{5. Áp dụng chính sách};
  \draw[msg] (9.7,3.0) -- (12.55,3.0) node[midway, above, align=center]{6. Video: còn hình,\\tắt tiếng};
  \draw[msg] (9.7,2.45) -- (12.55,2.45) node[midway, above, align=center]{7. Nhạc: phát tiếp,\\nguồn tiếng chính};
  \node[note] at (9.35,1.8) {tick ``Không hỏi lại'' $\Rightarrow$ lần sau áp dụng im lặng};
  \draw[dashed, black!55] (6.8,2.1) -- (7.4,1.86);
\end{tikzpicture}
\caption{Sequence diagram kịch bản ``chưa có quy tắc ghi nhớ'' của X0. Lựa chọn còn lại ``Dừng video'' dừng video và lưu vị trí vào Xem sau.}
\label{fig:seq-conflict}
\end{figure}

\subsubsection{Khởi động và khôi phục phiên}
\label{sec:seq-restore}

Hình~\ref{fig:seq-restore} mô tả những gì thực sự được khôi phục sau khi tiến trình bị hệ điều hành hủy, đối chiếu với mã nguồn hiện tại: tab cuối, bài nhạc cuối ở trạng thái dừng, danh sách phát, điểm và chuỗi ngày được nạp lại; vị trí video chỉ được phục hồi khi người dùng mở lại đúng video đó nhờ checkpoint trong Room (ghi mỗi 3 giây). Nhạc không tự phát lại và vị trí đọc tin chưa được lưu qua process death; hai điểm này nằm trong danh sách kế hoạch thay vì được mô tả như đã có.

\begin{figure}[H]
\centering
\begin{tikzpicture}[
  font=\tiny,
  ll/.style={draw=black!70, thick, rounded corners=2pt, fill=black!5, minimum width=2.15cm, minimum height=0.5cm, align=center},
  msg/.style={-{Latex[width=2mm, length=2.4mm]}, thick},
  ret/.style={-{Latex[width=2mm, length=2.4mm]}, thick, dashed},
  note/.style={draw=black!55, fill=black!3, rounded corners=2pt, inner sep=2.5pt, align=center}]
  \node[ll] (u)  at (1.1,6.45) {\bfseries Người dùng};
  \node[ll] (os) at (3.9,6.45) {Hệ điều hành};
  \node[ll] (app) at (6.8,6.45) {App\\(Splash $\rightarrow$ Main)};
  \node[ll] (st) at (9.7,6.45) {DataStore\\+ Room};
  \node[ll] (ss) at (12.55,6.45) {Session\\Coordinator};
  \foreach \n in {u,os,app,st,ss} \draw[dashed, black!50] (\n.south) -- ++(0,-4.6);
  % Thanh kích hoạt
  \foreach \bx/\ya/\yb in {6.8/3.9/5.4, 9.7/4.4/4.9, 12.55/3.1/3.9}{
    \draw[fill=white, draw=black!60, thick] (\bx-0.07,\ya) rectangle (\bx+0.07,\yb);}
  \draw[msg] (1.1,5.9) -- (3.9,5.9) node[midway, above]{1. Mở lại app (sau khi tiến trình bị hủy)};
  \draw[msg] (3.9,5.4) -- (6.8,5.4) node[midway, above]{2. Khởi động process + MainActivity};
  \draw[msg] (6.8,4.9) -- (9.7,4.9) node[midway, above]{3. Đọc phiên đã lưu, nạp danh sách phát};
  \draw[ret] (9.7,4.4) -- (6.8,4.4) node[midway, above]{4. tab cuối, bài cuối, playlists, điểm/streak};
  \draw[msg] (6.8,3.9) -- (12.55,3.9) node[midway, above]{5. restoreColdStartSession(): mở tab cuối; đặt bài cuối ở 0:00};
  \node[note] at (10.85,3.02) {nhạc không tự phát\\--~bấm play trên dock};
  \draw[dashed, black!55] (12.55,3.35) -- (12.1,3.08);
  \node[note, align=center] at (7.3,1.95) {Mở lại video dở: checkpoint Room 3 giây $\rightarrow$ seek đúng vị trí};
  \draw[dashed, black!55] (9.7,2.2) -- (9.55,2.02);
\end{tikzpicture}
\caption{Sequence diagram khôi phục phiên theo hành vi hiện thực. Vị trí đọc tin, ván game dở qua process death và tự phát lại nhạc thuộc hạng mục kế hoạch.}
\label{fig:seq-restore}
\end{figure}

\subsubsection{Xem video và chuyển phân hệ (X1)}
\label{sec:seq-pip}

Hình~\ref{fig:seq-pip} mô tả chuỗi quyết định then chốt của trải nghiệm đa nhiệm: trình phát đọc vị trí đã lưu trước khi mở stream; khi người dùng đổi tab, shell tính điều kiện \texttt{videoDocked} và hiển thị cửa sổ nổi dùng lại đúng engine thay vì tạo trình phát thứ hai; thao tác ``Mở rộng'' điều hướng về trình phát gốc nên vị trí phát liên tục. Tab Home là ngoại lệ: video phát ngay trong card của feed. Đóng cửa sổ nổi dừng video, lưu vị trí và thêm vào Xem sau.

\begin{figure}[H]
\centering
\begin{tikzpicture}[
  font=\tiny,
  ll/.style={draw=black!70, thick, rounded corners=2pt, fill=black!5, minimum width=2.15cm, minimum height=0.5cm, align=center},
  msg/.style={-{Latex[width=2mm, length=2.4mm]}, thick},
  ret/.style={-{Latex[width=2mm, length=2.4mm]}, thick, dashed},
  note/.style={draw=black!55, fill=black!3, rounded corners=2pt, inner sep=2.5pt, align=center}]
  \node[ll] (u)  at (1.1,6.45) {\bfseries Người dùng};
  \node[ll] (ui) at (3.9,6.45) {Màn Video (UI)};
  \node[ll] (sc) at (6.8,6.45) {Session\\Coordinator};
  \node[ll] (pv) at (9.7,6.45) {Video engine\\(Media3)};
  \node[ll] (sh) at (12.55,6.45) {Shell\\(cửa sổ nổi)};
  \foreach \n in {u,ui,sc,pv,sh} \draw[dashed, black!50] (\n.south) -- ++(0,-4.6);
  % Thanh kích hoạt
  \foreach \bx/\ya/\yb in {3.9/4.35/5.9, 6.8/4.85/5.4, 9.7/4.35/4.95, 12.55/2.5/3.95}{
    \draw[fill=white, draw=black!60, thick] (\bx-0.07,\ya) rectangle (\bx+0.07,\yb);}
  \draw[msg] (1.1,5.9) -- (3.9,5.9) node[midway, above]{1. Chạm video trong feed};
  \draw[msg] (3.9,5.4) -- (6.8,5.4) node[midway, above, align=center]{2. playVideo(id): đọc vị trí đã lưu; nhạc đang phát $\rightarrow$ popup C2 (X0)};
  \draw[msg] (6.8,4.9) -- (9.7,4.9) node[midway, above]{3. Mở stream tại resumeFrom};
  \draw[ret] (9.7,4.4) -- (3.9,4.4) node[midway, above]{4. Phát tiếp tại vị trí đã lưu};
  \draw[msg] (1.1,3.9) -- (12.55,3.9) node[midway, above]{5. Chọn tab khác (ví dụ Tin tức)};
  \node[note] at (10.8,3.05) {videoDocked = true\\cửa sổ nổi hiện, vẫn phát};
  \draw[dashed, black!55] (12.55,3.5) -- (12.1,3.1);
  \draw[msg] (12.55,2.5) -- (3.9,2.5) node[midway, above]{6. Chạm ``Mở rộng'': mở lại player, cùng engine, đúng vị trí};
  \node[note] at (9.2,1.75) {Đóng cửa sổ nổi: dừng + lưu vị trí + thêm Xem sau};
  \draw[dashed, black!55] (12.55,2.2) -- (10.85,1.8);
\end{tikzpicture}
\caption{Sequence diagram X1 --- cửa sổ nổi dùng chung engine với trình phát gốc; tab Home phát trong card feed thay vì mở cửa sổ nổi.}
\label{fig:seq-pip}
\end{figure}

\subsubsection{Tải danh sách tin theo offline-first}
\label{sec:seq-news-offline}

Hình~\ref{fig:seq-news-offline} mô tả luồng đọc dữ liệu chuẩn mà các phân hệ dùng chung, lấy tab Tin làm ví dụ. Room là nguồn sự thật duy nhất của UI: ViewModel chỉ quan sát Flow từ repository, nên bản cache hiện ngay (đáp ứng UX-01) trong khi request mạng chạy song song. Kết quả mạng không trả thẳng cho UI mà được ghi vào Room, và Flow tự phát bản mới. Khi mất mạng hoặc server lỗi, UI giữ nguyên bản cache; chỉ khi cache rỗng mới chuyển sang \texttt{EmptyState}/\texttt{ErrorState} kèm nút ``Thử lại''.

\begin{figure}[H]
\centering
\begin{tikzpicture}[y=1.35cm,
  font=\tiny,
  ll/.style={draw=black!70, thick, rounded corners=2pt, fill=black!5, minimum width=1.9cm, minimum height=0.5cm, align=center},
  msg/.style={-{Latex[width=2mm, length=2.4mm]}, thick},
  ret/.style={-{Latex[width=2mm, length=2.4mm]}, thick, dashed},
  note/.style={draw=black!55, fill=black!3, rounded corners=2pt, inner sep=2.5pt, align=left}]
  \node[ll] (u)  at (0.9,7.0) {\bfseries Người dùng};
  \node[ll] (ui) at (3.4,7.0) {Màn Tin (UI)};
  \node[ll] (vm) at (5.9,7.0) {NewsViewModel};
  \node[ll] (rp) at (8.4,7.0) {NewsRepository};
  \node[ll] (db) at (10.9,7.0) {Room};
  \node[ll] (api) at (13.4,7.0) {API server};
  \foreach \n in {u,ui,vm,rp,db,api} \draw[dashed, black!50] (\n.south) -- ++(0,-5.9);
  \draw[msg] (0.9,6.35) -- (3.4,6.35) node[midway, above]{1. Mở tab Tin};
  \draw[msg] (3.4,5.95) -- (5.9,5.95) node[midway, above]{2. collect uiState};
  \draw[msg] (5.9,5.55) -- (8.4,5.55) node[midway, above]{3. observeCatalog(): Flow};
  \draw[msg] (8.4,5.15) -- (10.9,5.15) node[midway, above]{4. truy vấn cache};
  \draw[ret] (10.9,4.75) -- (5.9,4.75) node[midway, above]{5. bản cache $\rightarrow$ UiState.Content (Loading nếu rỗng)};
  \draw[msg] (8.4,4.3) -- (13.4,4.3) node[midway, above]{6. GET /news/articles?page=1};
  % alt frame
  \draw[black!60, thick] (7.9,3.95) rectangle (14.1,1.25);
  \node[draw=black!60, fill=white, font=\tiny\bfseries, anchor=north west, inner sep=1.5pt] at (7.9,3.95) {alt};
  \node[font=\tiny, anchor=north west] at (8.45,3.95) {[200 OK]};
  \draw[ret] (13.4,3.5) -- (8.4,3.5) node[midway, above]{7. items, hasMore};
  \draw[msg] (8.4,3.1) -- (10.9,3.1) node[midway, above]{8. upsert vào Room};
  \draw[ret] (10.9,2.7) -- (5.9,2.7) node[midway, above]{9. Flow phát bản mới $\rightarrow$ UI cập nhật};
  \draw[black!60, dashed] (7.9,2.3) -- (14.1,2.3);
  \node[font=\tiny, anchor=north west] at (8.45,2.3) {[mất mạng / 5xx]};
  \draw[ret] (8.4,1.6) -- (5.9,1.6) node[midway, above]{10. lỗi thật};
  \node[note] at (2.8,1.6) {Giữ bản cache; cache rỗng $\rightarrow$\\EmptyState / ErrorState + ``Thử lại''};
\end{tikzpicture}
\caption{Sequence diagram tải danh sách tin theo offline-first: Room là nguồn sự thật của UI; kết quả mạng được ghi vào Room rồi mới tới UI.}
\label{fig:seq-news-offline}
\end{figure}

Ở bản nguyên mẫu, \texttt{NewsRepository} đọc danh mục mock và bài đã lưu nằm trong DataStore; bước gọi API server (6--10) là thiết kế giai đoạn remote. Do repository giữ nguyên contract \texttt{observeCatalog(): Flow}, việc bổ sung client Retrofit không làm thay đổi ViewModel hay màn hình.
